\section{Controller and Training Details}
\label{app:method}
\subsection{Consolidation configuration}
Algorithm~\ref{alg:engram} specifies the general multi-round controller. The scaled SFT experiments use a single batch update followed by per-skill verification and deployment; replay from earlier consolidation rounds is empty in this setting. Historical multi-round runs and their configurations are reported separately.

\subsection{Trajectory construction}
Acquisition targets low-mastery skills. A strong teacher executes action sequences through the same environment interface. The environment's reported target outcome determines demonstration success. For achievement-based environments, the stored segment ends at the action that unlocks the achievement, and a target-directed episode can stop once that event occurs. This alignment retains the action that caused success rather than an arbitrary prefix of the rollout.

Teacher demonstrations complement student successes and teacher-written failure corrections. A correction enters action training only if it supplies a valid action for an aligned observation. Source tags distinguish demonstrations, corrections, and student traces.

\subsection{Training configuration}
Examples from admitted skills are combined with available training-only replay:
\begin{equation}
 \dataset_t=\bigcup_{s:u_{t,s}=1}\dataset_{t,s},\qquad
 \widetilde{\dataset}_t=\dataset_t\cup
 \operatorname{Sample}(\replay_t,\rho).
\label{eq:dataset}
\end{equation}

The replay fraction $\rho$ is limited by archive availability. Each example contains task input, training context, and the target action. The scaled SFT runs include ETG context and optimize the action-token likelihood in Eq.~\eqref{eq:loss}; their candidates restart from the base model. Optional regularization adds
\begin{equation}
\beta\,\mathbb{E}_{(x,c,a)\sim\widetilde{\dataset}_t}
\operatorname{KL}\!\left(\pi_{\theta_0}(\cdot\mid x,c)\,\|\,\pi_\theta(\cdot\mid x,c)\right)
\end{equation}
at supervised response positions. This optional term is not part of the scaled SFT objective. Candidate loading precedes verification, and failures preserve existing approvals and external records.

\paragraph{Admission and probe settings.}
The demonstration threshold $d_{\min}=2$ counts successful records, not independent teacher episodes. QA probes use an answer-score threshold of 0.5; final EM and F1 are reported separately. Probe construction and episode separation follow Appendix~\ref{app:protocol}. Only records with valid training participation are eligible for release, with up to $k_{\mathrm{res}}$ recent residual records retained per memory category.

\paragraph{Retrieval-store accounting.}
If $A_t$ records arrive and $E_t$ eligible records are released, retrieval occupancy obeys
\begin{equation}
|\memory_{t+1}|=|\memory_t|+A_t-E_t.
\label{eq:occupancy}
\end{equation}
Retrieval storage stabilizes when departures keep pace with arrivals. Replay and any retained archives are accounted for separately; removing retrieval records does not by itself imply a reduction in total storage.
